AIsafePy is implemented in pure Python (3.10+) with approximately
5,000 lines of source code across 38 modules. The required runtime
dependencies are small: \texttt{pydantic}, \texttt{anyio},
\texttt{opentelemetry-api}, \texttt{tenacity}, and
\texttt{typing-extensions}. Heavier dependencies (Transformers, Torch,
HDBSCAN, sentence-transformers, scikit-learn) live behind optional
extras (\texttt{[stream]}, \texttt{[probes]}, \texttt{[adapt]}) so a
deployment in a constrained environment (such as a 512MB AWS Lambda
function) can use the deterministic regex and IFC paths without any
ML dependencies.

\paragraph{Public API surface.} The library exposes only what is
needed at each call site. \texttt{aisafepy.core} re-exports the
shared types; submodules are imported lazily on first attribute
access (Python's \texttt{\_\_getattr\_\_} module hook) so
\texttt{import aisafepy} is cheap even on machines where Torch is
installed. This matters in serverless cold-start scenarios.

\paragraph{Async-first.} Guards conform to the
\texttt{Callable[[Context], Awaitable[GuardDecision]]} protocol. The
pipeline runs guards within a tier in parallel via
\texttt{asyncio.gather}. Synchronous callers use \texttt{run\_sync},
which detects an already-running event loop and offloads to a worker
thread to avoid deadlocking under pytest-asyncio or Jupyter.

\paragraph{Adapter pattern for frameworks.} Each agent-framework
adapter is around 200 lines and follows the same shape: locate the
tool's callable, wrap it with a check-then-call proxy, register a
lifecycle hook for IFC context propagation, and attach a sentinel
attribute (\texttt{\_\_aisafepy\_policy\_\_}) so subsequent calls can
inspect the binding. The MCP adapter additionally patches the
server-side \texttt{add\_tool} and \texttt{handle\_call\_tool}
methods.

\paragraph{Testing.} The unit test suite has 59 tests covering core
primitives, IFC enforcement (including the AgentDojo
``important-message'' attack pattern), deterministic regex guards,
the cascading pipeline, adapter clustering, and the audit-log
hash-chain invariant. The full test suite runs in $<$1 second on
commodity hardware. Continuous integration runs on Python 3.10, 3.11,
and 3.12 in parallel.

\paragraph{Linting and reproducibility.} The repository uses
\texttt{ruff} for linting and formatting and pins a deterministic
toolchain in \texttt{pyproject.toml}. The
\texttt{[tool.hatch.build.targets.sdist]} section explicitly enumerates
the directories shipped in the sdist (source, tests, examples,
policies, docs) so downstream auditors can reproduce the artifact
from the GitHub tag.

\paragraph{Documentation and examples.} Four runnable examples
exercise each module end to end:
\texttt{flow\_\allowbreak openai\_\allowbreak agents.py}
demonstrates the indirect-injection attack and the IFC block;
\texttt{stream\_\allowbreak cascade.py} runs a three-tier
pipeline on a simulated streamed response;
\texttt{adapt\_\allowbreak pyrit\_\allowbreak to\_\allowbreak guards.py}
compiles a fabricated PyRIT-style log into runnable artifacts; and
\texttt{loop\_\allowbreak detection.py} demonstrates progress-tracker
triggers.

\paragraph{Distribution.} The library is published on PyPI under
\texttt{aisafepy}. The default \texttt{pip install aisafepy}
installs only the core. Extras follow the standard
\texttt{aisafepy[stream]} pattern.
